\documentclass[journal]{IEEEtran}
\usepackage{amsmath,amssymb,amsthm}
\usepackage{booktabs}
\newtheorem{theorem}{Theorem}
\newtheorem{theorem*}{Theorem}
\newtheorem{lemma*}{Lemma}
\newtheorem{proposition*}{Proposition}
\newtheorem{lemma}{Lemma}
\newtheorem{proposition}{Proposition}
\newtheorem{definition}{Definition}
\newcommand{\homc}{\mathrm{hom}}

\begin{document}
\title{Supplementary Material:\\
Proofs for ``Learnable and Transferable Logical Operators for Complex
Query Answering on Knowledge Graph Foundation Models''}

\author{}
\maketitle

\section{Notation and Conventions}

A knowledge graph is a finite structure $G=(V,R,E)$ with
$E\subseteq V\times R\times V$. $G^{+}$ denotes the closure under inverse
relations $\bar r$ (one fresh symbol per $r\in R$); all queries are
evaluated over $G^{+}$ so that atom directions are without loss of
generality. For a pattern (witness) structure $W$ with a distinguished
anchor node $\eta$ and $e\in V$, $\homc_W(G,e)$ is the number of
homomorphisms $W\to G^{+}$ with $\eta\mapsto e$. Queries follow
Definition~1 of the main text (EPFO with atom-level negation, DNF
disjunction); $\mathbf 1[\cdot]$ is the indicator.

\section{Proof of Theorem 1 (Canonical Representation)}

\begin{theorem*}
For every EPFO$^\neg$ query $q$ there is a finite family $\mathcal W_q$ of
anchored witness structures such that for every $G$ and $e\in V(G)$,
$\mathbf 1[e\in q(G)]$ is an integer combination of
$\{\homc_W(G,e)\}_{W\in\mathcal W_q}$; negated atoms enter through
inclusion--exclusion.
\end{theorem*}
\begin{proof}
\emph{Step 1 (positive conjunctive core).} Let $q^\circ$ be the
positive conjunctive part with variables $x_1,\dots,x_m$ (including the
answer variable $x_?$). Associated with $q^\circ$ is its \emph{canonical
database} $D(q^\circ)$: the typed pattern graph whose nodes are the
variables and whose typed edges are the atoms. By the classical
query-to-homomorphism correspondence,
$e\in q^\circ(G)\iff\exists$ homomorphism $h:D(q^\circ)\to G^{+}$ with
$h(x_?)=e$. Partitioning homomorphisms by the tuple of images
$(h(x_1),\dots,h(x_m))$ gives
\[
\homc_{D(q^\circ)}(G,e) \;=\; \sum_{(e_1,\dots,e_m):\,e_{?}=e}
\prod_{\text{atoms }(x_i,r,x_j)} \mu_G(e_i,r,e_j),
\]
where $\mu_G(a,r,b)=1$ if $(a,r,b)\in E(G^{+})$ and $0$ otherwise. Because
$D(q^\circ)$ encodes exactly the atoms, every summand is a product of
indicators of atoms; hence
\[
\mathbf 1[e\in q^\circ(G)] \;=\; \mathbf 1\!\left[\homc_{D(q^\circ)}(G,e) \ge 1\right].
\]
Set $\mathcal W_{q^\circ}=\{D(q^\circ)\}$ (a single witness of order
$m$); membership is decided by one anchored hom-count.

\emph{Step 2 (disjunction, DNF).} If $q = \bigvee_{\ell} c_\ell$ with
conjunctive clauses $c_\ell$, then
$\mathbf 1[e\in q(G)] = 1-\prod_\ell(1-\mathbf 1[e\in c_\ell(G)])$, a
polynomial in the clause indicators, each of which is Step-1 decidable.
The witness family is the union of the clause databases.

\emph{Step 3 (atom-level negation, inclusion--exclusion).} Let
$q = c \wedge \neg a_1 \wedge \cdots \wedge \neg a_k$ where $c$ is
conjunctive and the $a_i$ are negated atoms (possibly with fresh
existential variables, as in \texttt{pin}/\texttt{pni}). Then
\[
\mathbf 1[e\in q(G)] = \mathbf 1[e\in c(G)]\cdot
\prod_{i=1}^{k}\bigl(1-\mathbf 1[e\in (c\wedge a_i)(G)]\bigr).
\]
Expanding the product,
\[
\mathbf 1[e\in q(G)] = \sum_{S\subseteq[k]} (-1)^{|S|}\,
\mathbf 1\bigl[e\in (c\wedge\bigwedge_{i\in S}a_i)(G)\bigr],
\]
and each $(c\wedge\bigwedge_{i\in S}a_i)$ is again a positive conjunctive
query whose canonical database $D_S$ is a witness of order at most
$m+\sum_{i}|{\rm var}(a_i)|$. Hence
\[
\mathbf 1[e\in q(G)] \;=\; \sum_{S\subseteq[k]}(-1)^{|S|}\,
\mathbf 1\bigl[\homc_{D_S}(G,e)\ge 1\bigr],
\]
a finite integer combination of anchored hom-counts; negated atoms
contribute through inclusion--exclusion as claimed. Combining Steps 2 and
3 handles arbitrary DNF with negation. The bound on witness order is
$|{\rm var}(q)|+\sum_i |{\rm var}(a_i)|$, finite and independent of $G$.
\end{proof}

\emph{Remark.} The theorem states a sufficient statistic for
\emph{decidability}; executors in practice approximate thresholded counts
by smooth scores. The smooth statement used in the main text---order-$k$
executors factor through order-$k$ signatures---follows because every
feature consumed by such an executor is a function of the signature by
Definition~3, and signatures of order $k$ determine signatures of order
$\le k$.

\section{Proof of Theorem 2 (Separation)}

\begin{theorem*}
There is an explicit, linear-time constructible pair of query worlds
$(G_1,q_1)$, $(G_2,q_2)$ such that (i) every score-factored fuzzy pipeline
outputs identical rankings for $q_1$ in $G_1$ and $q_2$ in $G_2$, while
(ii) their ground truths differ and an order-2 witness aggregator
separates them.
\end{theorem*}
\begin{proof}
\emph{Construction.} Entities: anchors $a$, targets
$t_1,t_2,t_3,t_4$, marker hubs $z$. Relations: $r_a$; ambiguous
relations $r_1,r_2$; type markers $\tau_1,\tau_2$. In each \emph{block}
$b$: (a) $r_a$-edges from the anchor $a_b$ to all four targets;
(b) a hidden half $H_b\subset\{t_{b1},...,t_{b4}\}$ with $|H_b|=2$, chosen
uniformly; (c) each target in the $r_1$-half carries a $\tau_1$-edge and
each target in the $r_2$-half a $\tau_2$-edge; (d) in \emph{demonstrator}
blocks only, the local $r_1/r_2$ edges from the anchor are observed. The
query in a hidden block is $q_b=(a_b,r_a)\wedge(a_b,r_q)$ with
$r_q\in\{r_1,r_2\}$; its ground truth is the two targets of the
$r_q$-half, determined by the hidden edges---equivalently, by which
targets carry $\tau_{r_q}$-markers, since the construction ties the
marker type to the half.

\emph{(i) Score-factored pipelines collide.} Fix any executor whose
projection returns per-entity scores and whose logical operators are
pointwise functions of per-entity branch scores (products, probabilistic
sums, complements, or any fixed continuous combinations thereof). In a
hidden block the anchor has \emph{no observed} $r_q$-edges; conditional
on the block's local observations, the branch distribution over entities
is a function of $(a_b, r_q)$'s local features only. By construction the
local features of $(a_b,r_q)$ are identical across the two worlds (both
are ``anchor with $r_a$-fanout four and unobserved $r_q$''), and they are
identical for $r_q\in\{r_1,r_2\}$ within a world up to relabeling of
relation symbols, which transferability (Definition~2) forces to act as an
isomorphism. Hence the computed branch marginals, and therefore the final
pointwise combinations and rankings, coincide across $(G_1,q_1)$ and
$(G_2,q_2)$: the pipeline cannot be correct on both, because the ground
truths differ (disjoint answer sets). Its expected MRR equals that of
ranking four candidates of which two are relevant uniformly at random:
$\tfrac14(1+\tfrac12+\tfrac13+\tfrac14)\cdot 2=\tfrac{2}{3}$ rounded to
$0.67$ (matches the empirical $0.673$--$0.703$).

\emph{(ii) Order-2 aggregation separates.} Define the aggregator that
scores a target $t$ in branch $r_q$ by the indicator that $t$ carries a
$\tau_{r_q}$-edge. This is a function of the order-2 signature: the pair
$(r_q,\tau_{r_q})$ co-occurrence is a two-node pattern count in the
relation-interaction graph, and target markerhood is a one-hop anchored
count. The aggregator's intersection then returns exactly the two
relevant targets; MRR $=1$. This is precisely the rule our order-2 gate
variant implements (branch relation feature concatenated into the gate),
and the empirical $1.000$ matches.
\end{proof}

\emph{Remark (harmlessness in distribution).} When the projection is
trained to near-calibration on the query distribution, equal marginals
co-occur with equal groundings, so the collision of (i) does not reduce
accuracy---consistent with our paired-collision analysis on FB15k-237
(output Spearman $0.999$, MRR $\approx 1$). The separation bites exactly
when local evidence is missing or shifted, which is the zero-shot regime
of \S6.3.

\section{Proof of Theorem 3 (Learnable-Operator Characterization)}

\begin{theorem*}
The class $\mathcal O_k$ of order-$k$ operators satisfying transferability
is exactly the class of branch-permutation-invariant functions computable
from multisets of order-$k$ branch signatures. Fixed continuous
$t$-norms/$t$-conorms act pointwise on order-0 evidence and form a strict
subclass of measure zero inside any smoothly parameterized order-$\ge1$
family containing them.
\end{theorem*}
\begin{proof}
(\emph{Structure of the class.}) An operator $\omega$ consumes the branch
outputs of a plan prefix. Transferability (Definition~2) implies $\omega$
is invariant under entity and relation relabelings; since the plan's
branch list is unordered (conjunction and disjunction are commutative in
the query DAG), $\omega$ is additionally invariant under branch
permutations. By definition of aggregation order, $\omega$ is a measurable
function of the multiset $\{s_i,\sigma_i\}_{i\le m}$ of branch scores and
order-$k$ signatures. Conversely, any branch-symmetric measurable function
of those multisets defines an element of $\mathcal O_k$: it consumes no
identifiers, and multiset arguments make the relabeling action an
isomorphism, so equivariance holds.

(\emph{Fixed $t$-norms are order-0 and strict.}) A continuous $t$-norm
combination applied pointwise over entities depends on the branch score
vectors only---order-0 evidence---hence lands in
$\mathcal O_0\subsetneq\mathcal O_k$ for $k\ge1$; strictness is witnessed
by Theorem~2's separator, which is not expressible by any pointwise
combination.

(\emph{Measure zero.}) Let $\{T_\theta\}_{\theta\in\Theta}$ be a smoothly
parameterized family containing the product at $\theta^\ast$, e.g., the
Acz\'el--Alsina family with $\theta^\ast=1$, and let the gate be a smooth
function of parameters and evidence. The set of parameters for which the
operator equals $T_{\theta^\ast}$ pointwise is the preimage of a point
under a smooth map into a space of higher dimension; by Sard's theorem it
has measure zero. (For the one-parameter Acz\'el--Alsina family the
preimage of $\{\theta^\ast\}$ under a non-degenerate gate is a
codimension-one submanifold.)
\end{proof}

\section{Lemma 1 (Exact Degeneration)}

\begin{lemma*}
Setting every \textsc{Sqa}{} gate output to the constant $c$ with
$\sigma(c)=\tfrac{1-\theta_0}{3.5}\big|_{\theta_0=1}$, i.e.
$\theta\equiv1$, makes the gated intersection equal to the product
$t$-norm for all inputs, and dually for union.
\end{lemma*}
\begin{proof}
With $\theta_i\equiv 1$ for all branches,
\[
T_\Theta(s_1,\dots,s_m)=\exp\Big[-\sum_i(-\ln s_i)\Big]
=\prod_i \exp(\ln s_i)=\prod_i s_i,
\]
using $\sum_i(-\ln s_i)=-\ln\prod_i s_i$ and continuity of $\exp$; the
dual statement for union follows by applying the identity to complements
$1-s_i$. Numerically we verify $\max|T_\Theta-\prod s_i|<10^{-3}$ at the
forced gate on random inputs. Note the exponent $1/\theta_1=1$ as well,
so no branch asymmetry survives.
\end{proof}

\section{Proposition (Transfer Continuity)}

\begin{proposition*}
Define $d_{\mathrm{hom}}(G,G')$ as the $\ell_\infty$ discrepancy of
anchored witness-count profiles over all patterns of order $\le k$. For
any order-$k$ executor whose aggregation is $L$-Lipschitz in its inputs,
$|F(G,q)(e)-F(G',q)(e')|\le L\,d_{\mathrm{hom}}(G,G')$ for anchors $e,e'$
corresponding under the profile.
\end{proposition*}
\begin{proof}[Proof sketch]
The projection's messages are continuous functions of the anchored
profiles (each hop composes counts along observed edges, Lipschitzly in
the counts by the mean-value bound and the bounded degree of the plan).
The order-$k$ operator is $L$-Lipschitz in the branch scores and
signatures by hypothesis. Chaining the two Lipschitz maps and taking the
supremum over the finitely many order-$\le k$ patterns yields the claim.
The full multi-relational graphon formalism (existence of limits of
anchored profiles, tightness) follows the classical bounded-degree
convergence theory; we include it for completeness in the released
technical note.
\end{proof}

\section{Experimental Details Referred to in the Main Text}

\subsection{Hyperparameters}
Projection: 2-layer conditional propagation per hop, hidden 64.
Relation encoder: 3-layer typed MP over $G_R$, hidden 64. Gates: 2-layer
MLP ($65\!\to\!64\!\to\!1$; order-2 variant $129\!\to\!64\!\to\!1$),
$\theta\in(0.5,4)$ via $0.5+3.5\sigma(\cdot)$. Training: AdamW, ranking
loss (pretraining, lr $10^{-3}$, 50k steps, batch 12) or official
balanced logit-BCE (fine-tuning, lr $5\times10^{-4}$, 10k steps, batch
32). Total executor parameters $\approx 145$k (gated).

\subsection{Inference cost}
Measured on one RTX PRO 6000, FB15k-237 ($|V|{=}14{,}505$, $|E|{=}544$k
with inverses): fixed operators 21.3 it/s (batch 12, all types);
gated 19.6 it/s; gated order-2 19.1 it/s---an $8$--$10\%$ overhead over
fixed operators, negligible against projection cost. Zero-shot executor
fusion (router) doubles projection cost (two executors), $\sim$2$\times$.

\subsection{Gate statistics after official fine-tuning}
Empirical $\theta$ on sampled \texttt{2i} queries concentrates in
$[0.52,0.57]$---systematically \emph{softer} than the product point
$\theta=1$: the learned optimum relaxes conjunction strength rather than
sharpening it, consistent with the large disjunction gains
(\texttt{2u} $0.963$) under fuzzy margins.

\end{document}
